\documentclass[8pt,a4paper]{article}
\usepackage[utf8]{inputenc}
\usepackage[spanish]{babel}
\usepackage{graphicx}
\usepackage{amsmath}
\usepackage{listings}
\usepackage{xcolor}
\usepackage{geometry}
\usepackage{hyperref}
\usepackage{multicol}

\geometry{margin= 2.5cm}

% Configuración para código Python
\lstset{
    language=Python,
    basicstyle=\ttfamily\tiny,
    keywordstyle=\color{blue},
    commentstyle=\color{green!60!black},
    stringstyle=\color{red},
    numbers=left,
    numberstyle=\tiny\color{gray},
    breaklines=true,
    frame=single,
    backgroundcolor=\color{gray!10}
}

\usepackage{fontspec}
\setmainfont{EB Garamond}[
    UprightFont = * Regular,
    ItalicFont = * Italic,
    BoldFont = * SemiBold,
    BoldItalicFont = * SemiBold Italic,
]


% --- Espaciado ---
\usepackage{setspace}
\setstretch{1.2}
\setlength{\parskip}{0.2em}


\begin{document}

\section{Graph Neural Networks}

Una \emph{Graph Neural Network} (GNN) es un tipo de red neuronal diseñada para datos
con estructura de grafo: nodos (por ejemplo, personas, átomos, ciudades) y enlaces
entre ellos (amistades, enlaces químicos, carreteras, etc.).

La idea clave es que cada nodo no se procesa aislado, sino teniendo en cuenta a sus
vecinos en el grafo.

\section{Intuición de funcionamiento}

El proceso típico por capas es:

\begin{enumerate}
    \item Cada nodo empieza con un vector de características (edad, tipo de átomo, etc.).
    \item En cada capa, el nodo recibe información resumida de sus vecinos
    (por ejemplo, un promedio de sus vectores).
    \item El nodo mezcla su propia información con la de sus vecinos y actualiza su vector.
\end{enumerate}

Tras varias capas, el vector de cada nodo captura tanto sus atributos como el contexto
de su vecindario en el grafo.

\section{Ventajas y desventajas}

\subsection*{Ventajas}
\begin{itemize}
    \item Aprovechan explícitamente la estructura de conexiones entre entidades.
    \item Funcionan en dominios muy distintos: redes sociales, química, recomendación, etc.
    \item No dependen del orden en que numeramos los nodos.
\end{itemize}

\subsection*{Desventajas}
\begin{itemize}
    \item Pueden ser costosas en grafos muy grandes (problemas de memoria y tiempo).
    \item Muchas capas pueden hacer que los nodos se parezcan demasiado entre sí
    (sobre-suavizado).
    \item El diseño e implementación son más complejos que en redes totalmente conectadas
    o convolucionales básicas.
\end{itemize}

\section{Conclusión}

Las GNNs extienden las redes neuronales al caso en que la estructura de relaciones
entre los datos es tan importante como las características individuales.
Son muy útiles cuando el ``con quién estás conectado'' importa tanto como ``quién eres''. Una GNN aprende representaciones de nodos usando su información y la de sus vecinos en un grafo.

\newpage

\section{¿Qué es una Red Bayesiana?}

Una \emph{Red Bayesiana} es un modelo probabilístico que representa un conjunto
de variables y sus dependencias mediante un grafo dirigido sin ciclos.

\begin{itemize}
    \item Cada nodo = una variable (por ejemplo, enfermedad, síntoma, lluvia).
    \item Cada flecha = relación de dependencia causal o probabilística.
    \item Cada nodo tiene una tabla de probabilidades condicionadas a sus ``padres''.
\end{itemize}

La idea central es combinar teoría de probabilidad y estructura de grafo para
razonar con incertidumbre.

\section{Intuición de funcionamiento}

\begin{enumerate}
    \item Se construye el grafo indicando qué variable influye en cuál.
    \item A cada nodo se le asignan probabilidades condicionales (por ejemplo,
    probabilidad de fiebre dado que hay gripe).
    \item Dado un conjunto de evidencias (variables observadas), se actualizan
    las probabilidades de las demás variables usando la regla de Bayes.
\end{enumerate}

Así, la red permite contestar preguntas del tipo: ``dado lo que observo,
¿cuán probable es cada hipótesis?''.

\section{Ventajas y desventajas}

\subsection*{Ventajas}
\begin{itemize}
    \item Modelo explícito y gráfico de dependencias (muy interpretable).
    \item Manejan bien la incertidumbre y datos faltantes.
    \item Permiten combinar conocimiento experto y datos para aprender parámetros.
\end{itemize}

\subsection*{Desventajas}
\begin{itemize}
    \item Diseñar la estructura del grafo puede ser difícil en problemas grandes.
    \item Inferencia exacta puede ser muy costosa en redes complejas.
    \item Aprender estructura automáticamente requiere muchos datos y es computacionalmente caro.
\end{itemize}

\section{Conclusión}

Las Redes Bayesianas son una herramienta potente para modelar y razonar con
incertidumbre cuando existen relaciones causales o de dependencia claras
entre variables. Son especialmente útiles cuando, además de predecir,
queremos \emph{explicar} y \emph{entender} el comportamiento del sistema. Una Red Bayesiana modela dependencias probabilísticas entre variables mediante un grafo dirigido.

\end{document}
